\chapter{Tecnologie utilizzate}\label{chap:second}
Per far fronte alle sfide di estrazione, elaborazione, archiviazione e interrogazione dei documenti normativi, è stato impiegato uno stack tecnologico che combina strumenti di parsing avanzato, modelli di LLM ed embedding e database specializzati. Di seguito vengono presentate le tecnologie adottate, illustrando per ciascuna le motivazioni che ne hanno guidato la scelta.

\section{Linguaggi di Programmazione}
\begin{itemize}
    \item \textbf{Python}: È stato il linguaggio cardine per la costruzione dell'intera pipeline del progetto. La scelta è ricaduta su Python per diverse ragioni:
    \begin{itemize}
        \item \textbf{Ecosistema dedicato ai modelli linguistici:} Python dispone di librerie mature e ampiamente supportate (come \texttt{sentence-transformers}) che hanno permesso di integrare agevolmente i modelli di embedding senza dover reimplementare componenti di basso livello.
        \item \textbf{Ottima integrazione con tutte le librerie utilizzate:} Python risulta perfettamente integrato con tutti gli strumenti e le librerie adottate nell'intero stack tecnologico (dai tool di estrazione testuale, ai modelli LLM, fino alla Graph Data Science Library e al database Neo4j tramite API native), permettendo di orchestrare l'intera pipeline in modo fluido e senza la necessità di complessi adattatori o wrapper di terze parti.
    \end{itemize}
\end{itemize}

\section{Strumenti per l'estrazione del Testo}
La fase di estrazione ha dovuto confrontarsi con documenti PDF eterogenei: alcuni nativamente digitali e ben strutturati, altri derivanti da scansioni cartacee di qualità variabile. Per questo motivo è stata adottata una combinazione di strumenti complementari.

\begin{itemize}
    \item \textbf{Open Data Loader}\cite{opendataloader}: Libreria open-source per l'estrazione e il caricamento di dati da fonti eterogenee. È stata selezionata per due caratteristiche fondamentali:
    \begin{itemize}
        \item \textbf{Determinismo:} A differenza di approcci basati su modelli probabilistici, Open Data Loader estrae il testo in modo deterministico, garantendo coerenza dei risultati, un requisito indispensabile nel dominio legale, dove ogni modifica del testo estratto tra un'esecuzione e l'altra minerebbe la fiducia nel sistema.
        \item \textbf{Conversione diretta in Markdown:} Il Markdown si rivela un formato testuale ottimale per i LLM. La sua sintassi leggera preserva fedelmente la struttura logica del documento originale, fornendo al modello un contesto già strutturato che riduce significativamente gli errori nella successiva fase di parsing gerarchico.
    \end{itemize}

    \item \textbf{Docling}\cite{docling}: Libreria specializzata nell'analisi del layout documentale e nell'estrazione di contenuti strutturati. È stata integrata come supporto complementare a Open Data Loader per la gestione di PDF più complessi: supporta documenti multi-colonna, tabelle e layout non lineari, frequenti negli atti della pubblica amministrazione.

    \item \textbf{RapidOCR}\cite{rapidocr}: Motore di \textit{Optical Character Recognition} impiegato per digitalizzare i documenti derivanti da scansioni visive. La scelta è stata guidata da un'analisi comparativa tra diversi motori OCR:
    \begin{enumerate}
        \item \textbf{EasyOCR:} Configurato di default all'interno di Docling, adotta modelli basati su Deep Learning altamente accurati ma estremamente esigenti in termini computazionali. Anche isolando l'elaborazione sulla singola pagina, il consumo di RAM era così elevato che rendeva impossibile l'utilizzo nell'ambiente di produzione containerizzato.
        \item \textbf{Tesseract OCR:} Si è dimostrato estremamente leggero e pienamente compatibile con i limiti di memoria imposti, scongiurando qualsiasi crash del container. Tuttavia, nel contesto specifico del linguaggio legale italiano, il motore non è stato in grado di interpretare correttamente il layout, fallendo l'estrazione del testo e restituendo output incompleti o disallineati rispetto alla struttura originale del documento.
        \item \textbf{RapidOCR:} Questo specifico modulo è stato infine selezionato per l'estrazione del testo dalle scansioni visive. RapidOCR si è rivelato il compromesso ideale, offrendo un perfetto bilanciamento tra l'accuratezza del riconoscimento del testo italiano e il consumo di RAM del container.
    \end{enumerate}
\end{itemize}

\section{Scelta del Database}
La duplice natura del problema — da un lato la necessità di rappresentare relazioni gerarchiche complesse, dall'altro l'esigenza di una ricerca semantica efficiente — ha portato all'adozione di un database a grafo specializzato.
La scelta è ricaduta su Neo4j, il database a grafo più diffuso e maturo sul mercato, per le seguenti motivazioni: 
    \begin{itemize}
        \item \textbf{Modellazione naturale della gerarchia normativa:} La struttura di un atto normativo è intrinsecamente un grafo, in cui ogni elemento è un nodo e le relazioni di appartenenza sono archi. Neo4j consente di mappare questa topologia in modo diretto, senza le complessità delle \texttt{JOIN} ricorsive tipiche dei database relazionali.
        \item \textbf{Supporto nativo per la Vector Search:} Il database permette di memorizzare rappresentazioni vettoriali come proprietà dei nodi e di creare indici dedicati, configurabili con funzioni di similarità coseno o euclidea. Nelle versioni recenti di Neo4j, tali indici supportano vettori con dimensionalità compresa tra 1 e 4096 \cite{neo4j_vector_indexes}. Questa funzionalità risulta fondamentale per le operazioni di ricerca basate sulla \textit{cosine similarity}.
    \end{itemize}

Per interagire con il database e gestire le operazioni di interrogazione e popolamento del grafo, è stato adottato \textbf{Cypher} \cite{cypher}, il linguaggio dichiarativo nativo dell'ecosistema Neo4j.

Infine, per valorizzare le informazioni topologiche della rete, si è fatto ricorso alla \textbf{Graph Data Science Library}, una libreria nativa dedicata all'implementazione di algoritmi analitici sui grafi. Al suo interno è stato selezionato \textbf{Node2Vec}, un potente algoritmo di \textit{graph embedding}.
A differenza degli embedding semantici utilizzati per la ricerca testuale (discussi nella Sezione \ref{sec:llm_embedding_models}), questo algoritmo permette di codificare in uno spazio vettoriale di 128 dimensioni le complesse relazioni gerarchiche e di vicinato che intercorrono tra i vari elementi del corpus normativo.

\section{LLM e Modelli di Embedding Semantici}\label{sec:llm_embedding_models}
Il cuore intelligente del sistema risiede nei modelli di linguaggio e di embedding, responsabili rispettivamente della strutturazione semantica dei documenti e della loro vettorializzazione per la ricerca.

\begin{itemize}
    \item \textbf{Gemini 3.5 Flash}\cite{gemini}: \textit{Large Language Model} di Google, utilizzato tramite API per il compito cruciale della strutturazione dei documenti normativi. Le motivazioni alla base di questa scelta sono:
    \begin{itemize}
        \item \textbf{Aderenza rigorosa allo schema JSON:} Gemini 3.5 Flash supporta la generazione di output strutturati conformi a uno schema JSON predefinito \cite{gemini_structured_output}. Nell'ambito del progetto, questa funzionalità ha consentito di rappresentare la gerarchia normativa (Titoli, Capi, Articoli, Commi) in modo consistente e automaticamente processabile dalle fasi successive della pipeline.
        \item \textbf{Comprensione della lingua italiana:} Il modello ha mostrato una solida competenza nell'elaborazione del linguaggio burocratico italiano, riconoscendo correttamente la terminologia giuridica e le espressioni formulari tipiche degli atti amministrativi.
        \item \textbf{Rapporto costo/qualità:} Rispetto a modelli più potenti ma significativamente più costosi (come GPT-4 o Claude), Gemini 3.5 Flash ha offerto un compromesso ottimale tra qualità dell'output strutturato e costo per token, aspetto rilevante considerando l'elaborazione batch di centinaia di documenti.
    \end{itemize}

    \item \textbf{Modelli di Embedding:} La vettorializzazione del testo normativo rappresenta un elemento critico per la qualità della ricerca semantica. Sono stati testati molteplici modelli per identificare la soluzione più efficace nel dominio giuridico italiano. I test più approfonditi hanno coinvolto i seguenti modelli:
    \begin{itemize}
        \item \textbf{\texttt{google/embeddinggemma-300m}:} Modello sviluppato da Google che trasforma il testo in vettori a 768 dimensioni. È costruito sull'architettura Gemma 3 ed è addestrato su oltre 100 lingue. Mediante la tecnica \textit{Matryoshka Representation Learning}, la rappresentazione può essere ridotta a 512, 256 o 128 dimensioni, introducendo un compromesso tra compattezza e qualità dell'embedding \cite{embeddinggemma}.
        
        \item \textbf{\texttt{perplexity-ai/pplx-embed-v1-0.6b}:} Modello sviluppato da Perplexity AI che produce vettori a 1024 dimensioni ed è progettato per compiti di retrieval multilingue, tra cui la ricerca semantica e i sistemi RAG \cite{eslami2026pplxembed}.
        
        \item \textbf{\texttt{perplexity-ai/pplx-embed-context-v1-0.6b}:} Variante contestuale del modello Perplexity, anch'essa produce vettori a 1024 dimensioni. Durante la rappresentazione dei singoli chunk incorpora informazioni provenienti dal contesto globale del documento, anziché trattare ciascun frammento in modo completamente isolato \cite{eslami2026pplxembed}. Il modello applica il principio del \textit{Late Chunking} \cite{latechunking}, con l'obiettivo di preservare nelle rappresentazioni locali una parte dell'informazione disponibile nel documento completo.
\end{itemize}
    La valutazione si è basata sulla capacità di ciascun modello di recuperare i commi pertinenti a partire da query in linguaggio naturale. A questo scopo, la \textit{Precision@k} è stata impiegata per misurare la quota di risultati pertinenti presente nelle prime $k$ posizioni della graduatoria, mentre il \textit{Mean Reciprocal Rank} (MRR) ha permesso di valutare quanto precocemente comparisse il primo risultato pertinente. Si tratta di due metriche consolidate nell'ambito dell'Information Retrieval, complementari perché descrivono rispettivamente la qualità complessiva dei primi risultati e la rapidità con cui il sistema individua una risposta rilevante \cite{manning2008introduction}.
\end{itemize}

\section{Deployment e Containerizzazione}
\begin{itemize}
    \item \textbf{Docker}: L'intero ecosistema è stato modularizzato attraverso lo sviluppo di container isolati e specifici per ciascun componente della pipeline. Questa scelta architetturale è stata dettata da esigenze sia tecniche che operative:
    \begin{itemize}
        \item \textbf{Isolamento delle dipendenze:} Ogni fase della pipeline (estrazione, strutturazione, ingestion, embedding) presenta dipendenze software potenzialmente conflittuali. Docker ha permesso di incapsulare ciascuna in un ambiente dedicato, eliminando i conflitti di versione tra librerie e garantendo la riproducibilità dell'ambiente di esecuzione su qualsiasi macchina.
        \item \textbf{Modularità e manutenibilità:} La separazione in container indipendenti ha permesso di sviluppare, testare e aggiornare ciascun modulo della pipeline in isolamento, semplificando la manutenzione.
        \item \textbf{Scalabilità orizzontale:} L'architettura a container costituisce un prerequisito naturale per futuri scenari di deployment distribuito in cui diverse istanze della pipeline possano operare concorrentemente su lotti distinti di documenti.
    \end{itemize}
\end{itemize}
